\documentclass[preprint]{article}
\usepackage{neurips_2025}
\usepackage[utf8]{inputenc}
\usepackage[T1]{fontenc}
\usepackage{amsmath,amssymb,amsfonts}
\usepackage{booktabs}
\usepackage{graphicx}
\usepackage{hyperref}
\usepackage{algorithm}
\usepackage{algorithmic}
\usepackage{multirow}
\usepackage{xcolor}
\usepackage{natbib}

% Custom commands
\newcommand{\SketchBudget}{\textsc{SketchBudget}}
\newcommand{\BF}{\textrm{BF}}
\newcommand{\CMS}{\textrm{CMS}}
\newcommand{\HLL}{\textrm{HLL}}
\newcommand{\fpr}{\textrm{FPR}}
\newcommand{\Real}{\mathbb{R}}

\title{Optimal Error Budgeting for Heterogeneous Sketch Pipelines\\in Approximate Stream Processing}

\author{
  Anonymous Author(s)
}

\begin{document}
\maketitle

% ============================================================
% ABSTRACT
% ============================================================
\begin{abstract}
Modern stream processing systems compose multiple probabilistic data structures---Bloom filters, Count-Min Sketches, HyperLogLog counters---into pipelines to answer complex queries over high-velocity data streams. Current practice treats each sketch independently, using standalone error bounds and ad-hoc memory allocation, which leads to loose end-to-end guarantees and wasted memory. We propose \SketchBudget{}, a framework for analyzing error propagation through heterogeneous sketch pipelines and optimally distributing a fixed memory budget across pipeline stages. We formalize three canonical composition patterns (Filter-Estimate, Estimate-Threshold-Count, Filter-Estimate-Aggregate), derive composition bounds tighter than naive approaches by exploiting cross-stage error structure, and solve the resulting constrained optimization problem. On synthetic Zipfian streams and real NASA HTTP access logs, \SketchBudget{} reduces end-to-end error by 42--51\% compared to the best baseline on two of three pipeline configurations at 500KB total budget, and achieves equivalent accuracy with up to 47--72\% less memory. Our analysis reveals that composition-aware allocation is most beneficial for pipelines where filtering errors propagate multiplicatively to downstream stages, while threshold-based pipelines benefit less due to the dominance of frequency-distribution effects near the decision boundary.
\end{abstract}

% ============================================================
% 1. INTRODUCTION
% ============================================================
\section{Introduction}

Probabilistic data structures---Bloom filters~\citep{bloom1970space}, Count-Min Sketches~\citep{Cormode2005CountMin}, HyperLogLog counters~\citep{flajolet2007hyperloglog}---are fundamental building blocks in modern stream processing. Each provides approximate answers to a specific query type (set membership, frequency estimation, cardinality counting) with tunable accuracy-space tradeoffs. In practice, these structures are routinely composed into \emph{pipelines}: a network monitoring system might use a Bloom filter to check whether an IP belongs to a watchlist, a Count-Min Sketch to estimate traffic frequency of matched IPs, and a HyperLogLog to count distinct heavy hitters.

When sketches are composed in a pipeline, their errors interact: a Bloom filter's false positives create ``phantom items'' that inflate downstream frequency estimates; a Count-Min Sketch's overestimates shift the effective threshold for heavy-hitter detection, altering which items a downstream HyperLogLog counts. Despite this, current practice treats each sketch independently: practitioners size each structure using its standalone worst-case bound, then compose results with conservative union bounds. This leads to three problems: (1)~\emph{loose bounds}---union bounds dramatically overestimate end-to-end error; (2)~\emph{suboptimal allocation}---without understanding error interactions, systems allocate memory uniformly or by ad-hoc heuristics; and (3)~\emph{heterogeneous error types}---different sketches produce fundamentally different error distributions (one-sided vs.\ symmetric, additive vs.\ multiplicative) that interact non-trivially.

We propose \SketchBudget{}, a framework with three components:
\begin{itemize}
    \item \textbf{Error Propagation Algebra}: A formal model for how false positive rates, additive errors, and multiplicative errors compose through pipeline operators (filter, threshold, aggregate).
    \item \textbf{Tighter Composition Bounds}: For canonical pipeline patterns, we derive closed-form end-to-end error expressions as functions of per-stage parameters that are provably tighter than naive bounds, though still conservative in absolute terms.
    \item \textbf{Optimal Allocation Algorithm}: Given a pipeline and total memory budget, we formulate and solve the constrained optimization of distributing memory across stages to minimize end-to-end error.
\end{itemize}

Our contributions are:
\begin{itemize}
    \item We formalize error propagation through three canonical heterogeneous sketch pipeline patterns---the first such treatment for sequential pipelines of mixed sketch types.
    \item We derive composition bounds that are provably tighter than naive bounds: 0.6--4.5\% tighter for Filter-Estimate, 2--42\% for Filter-Estimate-Aggregate, and orders of magnitude tighter for Estimate-Threshold-Count pipelines. We note that all bounds remain conservative (13--1537$\times$ overprediction at $\alpha=1.0$), indicating that the primary value of our analysis is in enabling better \emph{relative} allocation decisions rather than precise error prediction.
    \item We show that composition-aware allocation reduces end-to-end error by 45\% (P1) and 42\% (P3) over the best baseline at 500KB budget, while honestly characterizing the P2 configuration where the Proportional baseline remains competitive.
    \item We evaluate on both synthetic Zipfian streams and real NASA HTTP access logs, with ablations over budget, distribution skew, and optimizer comparison.
\end{itemize}

The remainder of this paper is organized as follows. Section~\ref{sec:related} reviews related work. Section~\ref{sec:method} describes our error propagation algebra, composition bounds, and allocation algorithm. Section~\ref{sec:experiments} presents experimental results. Section~\ref{sec:discussion} discusses limitations. Section~\ref{sec:conclusion} concludes.

% ============================================================
% 2. RELATED WORK
% ============================================================
\section{Related Work}
\label{sec:related}

\paragraph{Foundational Probabilistic Data Structures.}
Bloom filters~\citep{bloom1970space} provide space-efficient set membership testing with tunable false positive rates. The Count-Min Sketch~\citep{Cormode2005CountMin} enables frequency estimation with one-sided additive error guarantees. HyperLogLog~\citep{flajolet2007hyperloglog} achieves near-optimal cardinality estimation using $O(\log \log n)$ bits per register. Quantile summaries~\citep{greenwald2001space} provide streaming quantile approximation. Our work uses these as pipeline components and analyzes how their errors propagate through composition.

\paragraph{Multi-Query Sketch Optimization.}
\citet{dobra2004sketch} studied sketch-based multi-query optimization over data streams, showing that space allocation for average error across parallel AMS sketch queries is NP-complete. \citet{dobra2009multi} extended this to a journal treatment with join coalescing heuristics. Critically, both works consider \emph{parallel queries} over \emph{homogeneous sketches} (all AMS/CMS), while we consider \emph{sequential pipelines} of \emph{heterogeneous sketch types} with cross-stage error propagation---a fundamentally different optimization landscape.

\paragraph{Learned and Adaptive Sketches.}
The learned index structures framework~\citep{kraska2018case} sparked interest in ML-augmented data structures. \citet{mitzenmacher2018model} formalized learned Bloom filters with the ``sandwiching'' optimization for improved false positive rates. \citet{dolera2023learning} applied Bayesian nonparametric priors to Count-Min Sketches for improved single-stage estimation. \citet{zhu2025mnemosyne} proposed workload-aware Bloom filter tuning for LSM trees. These works optimize individual sketch components; our framework is complementary, as improved single-sketch error models can be plugged into our pipeline-level optimization.

\paragraph{Positioning.}
No prior work provides (a)~a formal model of error propagation through sequential pipelines of heterogeneous sketch types, (b)~composition bounds tighter than naive approaches by exploiting cross-stage independence, or (c)~an optimal allocation algorithm for heterogeneous pipelines. Our work fills this gap.

% ============================================================
% 3. METHOD
% ============================================================
\section{Method}
\label{sec:method}

\subsection{Problem Formulation}

Consider a sketch pipeline as a directed acyclic graph where each node $i$ is a sketch operator. Each node has: an \emph{error function} $e_i(m_i)$ mapping allocated memory $m_i$ to its standalone error parameter; an \emph{error type} (false positive rate, additive error, or relative error); and a \emph{propagation function} describing how upstream errors affect its effective error.

Given total memory budget $M$, the allocation problem is:
\begin{equation}
\label{eq:optimization}
\min_{m_1, \ldots, m_k} \; E_{\text{total}}(m_1, \ldots, m_k) \quad \text{s.t.} \quad \sum_{i=1}^{k} m_i \leq M, \quad m_i \geq m_i^{\min} \; \forall i
\end{equation}
where $E_{\text{total}}$ is the composed end-to-end error from our error algebra, and $m_i^{\min}$ is the minimum memory for each sketch type to function correctly.

\subsection{Canonical Pipeline Patterns}

We formalize three patterns that cover common streaming analytics use cases.

\paragraph{P1: Filter-Estimate (BF $\to$ CMS).}
A Bloom filter with false positive rate $p$ filters items before frequency estimation by a Count-Min Sketch with error parameter $\varepsilon$. For a query set of $n_+$ true positives and $n_-$ true negatives, the end-to-end mean absolute error is:
\begin{equation}
\label{eq:p1_naive}
E_{\text{naive}}^{P1} = \frac{n_+ \cdot \varepsilon N + p \cdot n_- \cdot \varepsilon N}{n_+ + n_-}
\end{equation}
where $N$ is the stream length. The naive bound treats false-positive items identically to true positives, assigning each the worst-case CMS error $\varepsilon N$.

Our tight bound observes that false-positive items are \emph{absent} from the stream, so their CMS estimate equals only the hash collision overestimate. The expected overestimate for an absent item is $\varepsilon N / e$ (the mean of the minimum of $d$ i.i.d.\ exponential-tailed collision counts), yielding:
\begin{equation}
\label{eq:p1_tight}
E_{\text{tight}}^{P1} = \frac{n_+ \cdot \varepsilon N + p \cdot n_- \cdot \varepsilon N / e}{n_+ + n_-}
\end{equation}
The improvement ratio is $1 - (1/e) \cdot (p \cdot n_-)/(n_+ + p \cdot n_-)$, which ranges from 0\% (when $p=0$, no false positives) to $(1-1/e) \approx 63\%$ (when false positives dominate). For typical operating points with small $p$, this yields modest improvements of 0.6--4.5\%.

\paragraph{P2: Estimate-Threshold-Count (CMS $\to$ Threshold $\to$ HLL).}
A CMS estimates item frequencies; items exceeding threshold $T$ are declared heavy hitters; an HLL counts distinct heavy hitters. The naive bound treats all $n_d$ distinct items as potential false heavy hitters:
\begin{equation}
E_{\text{naive}}^{P2} = n_d \cdot \min\!\left(1, \frac{\varepsilon N}{T}\right) + \sigma_{\HLL} \cdot n_d
\end{equation}
where $\sigma_{\HLL} = 1.04/\sqrt{m_{\HLL}}$ is the HLL standard error.

Our tight bound uses the actual frequency distribution. An item with true frequency $f < T$ is falsely declared heavy only when the CMS overestimate exceeds $T - f$. Using a linear approximation for the CMS error CDF:
\begin{equation}
\label{eq:p2_tight}
E_{\text{tight}}^{P2} = \sum_{\{i : f_i < T\}} \max\!\left(0, \frac{\varepsilon N - (T - f_i)}{\varepsilon N}\right) + \sigma_{\HLL} \cdot (|\text{true HH}| + \hat{F})
\end{equation}
where $\hat{F}$ is the expected number of false heavy hitters from the distribution-aware sum. This bound can be orders of magnitude tighter than the naive bound because most items have frequencies far below the threshold, making the per-item false heavy hitter probability negligible.

\paragraph{P3: Filter-Estimate-Aggregate (BF $\to$ CMS $\to$ Sum).}
A Bloom filter filters for set $S$ membership; a CMS estimates per-item frequencies; results are summed. The naive bound assumes false positives get the worst-case CMS estimate:
\begin{equation}
E_{\text{naive}}^{P3} = |S| \cdot \varepsilon N + p \cdot n_{\text{neg}} \cdot \varepsilon N
\end{equation}
The tight bound uses the average CMS estimate for absent items ($\varepsilon N / e$):
\begin{equation}
\label{eq:p3_tight}
E_{\text{tight}}^{P3} = |S| \cdot \varepsilon N + p \cdot n_{\text{neg}} \cdot \frac{\varepsilon N}{e}
\end{equation}
The false-positive term is reduced by a factor of $1/e \approx 0.37$, yielding 2--42\% tighter bounds depending on the relative magnitude of the two terms.

\subsection{Optimal Allocation Algorithm}

Given the composition bounds, we formulate memory allocation as the constrained optimization in Eq.~\ref{eq:optimization}. The error functions for individual sketches are:
\begin{itemize}
    \item \textbf{Bloom filter}: $\fpr(m) = (1 - e^{-k^* n / (8m)})^{k^*}$ where $k^* = (8m/n) \ln 2$ is the optimal number of hash functions, $n$ is the expected number of items, and $m$ is memory in bytes.
    \item \textbf{Count-Min Sketch}: $\varepsilon(m) = e / w$ where $w = m/(4d)$ is the table width for depth $d$ with 4-byte counters.
    \item \textbf{HyperLogLog}: $\sigma(m) = 1.04/\sqrt{m}$ for $m$ registers.
\end{itemize}

For two-stage pipelines (P1, P3), the composed error is a function of a single free variable (given the budget constraint $m_1 + m_2 = M$), and we solve it exactly using bounded scalar optimization (\texttt{scipy.optimize.minimize\_scalar}). For P2, we enumerate discrete HLL precisions (powers of 2) and select the CMS/HLL split that minimizes the composed error.

We also implement a greedy marginal heuristic: starting from minimum allocations, iteratively assign the next memory increment $\Delta m$ to the stage with the largest marginal error reduction. We compare this against the exact optimizer in our ablation study.

% ============================================================
% 4. EXPERIMENTS
% ============================================================
\section{Experiments}
\label{sec:experiments}

\subsection{Setup}

\paragraph{Implementation.}
All sketches are implemented from scratch in Python with NumPy/SciPy for precise memory control. Bloom filters use \texttt{mmh3} hashing; CMS and HLL use \texttt{xxhash}. All experiments run on CPU (2 cores, 128GB RAM).

\paragraph{Datasets.}
We evaluate on three data sources:
\begin{itemize}
    \item \textbf{Zipfian synthetic}: $10^6$ items drawn from Zipf($\alpha$) with $\alpha \in \{0.5, 0.8, 1.0, 1.2, 1.5, 2.0\}$ over a universe of 100K items. Default: $\alpha=1.0$.
    \item \textbf{Synthetic network trace}: Heavy-tailed IP traffic with Zipf($\alpha=1.1$), 100K IPs, $10^6$ packets, 1000-item watchlist.
    \item \textbf{NASA HTTP access logs}: Real web server logs from July 1995~\citep{nasa_http}. We extract hostnames as stream items and use the first $10^6$ entries, yielding ${\sim}51$K distinct hosts.
\end{itemize}

\paragraph{Pipeline Configurations.}
We evaluate three canonical patterns:
\begin{itemize}
    \item \textbf{P1} (Filter-Estimate): BF $\to$ CMS. Query: estimate frequency of items in set $S$.
    \item \textbf{P2} (Estimate-Threshold-Count): CMS $\to$ Threshold $\to$ HLL. Query: count distinct heavy hitters (items with frequency $\geq T$, where $T = 10\times$ average frequency).
    \item \textbf{P3} (Filter-Estimate-Aggregate): BF $\to$ CMS $\to$ Sum. Query: total traffic from items in watchlist.
\end{itemize}

\paragraph{Baselines.}
We compare four allocation strategies:
\begin{itemize}
    \item \textbf{Uniform}: Equal memory per stage.
    \item \textbf{Independent}: Memory proportional to each sketch's standalone space requirement for equal per-stage error.
    \item \textbf{Proportional}: Fixed 30:70 (BF:CMS) or 90:10 (CMS:HLL) splits based on typical space complexity ratios.
    \item \textbf{\SketchBudget{} (ours)}: Composition-aware optimal allocation.
\end{itemize}

\paragraph{Metrics.}
Primary metrics: mean absolute error (P1), cardinality error (P2), absolute aggregate error (P3). All results are mean $\pm$ std over 5 random seeds $\{42, 123, 456, 789, 1024\}$.

\subsection{Main Results}

Table~\ref{tab:main_results} shows end-to-end error at 500KB budget on Zipfian ($\alpha=1.0$) data.

\begin{table}[t]
\centering
\caption{End-to-end error (mean $\pm$ std) at 500KB budget on Zipfian ($\alpha=1.0$) data. Lower is better ($\downarrow$). Best result in \textbf{bold}. Improvement is \SketchBudget{} vs.\ best baseline.}
\label{tab:main_results}
\begin{tabular}{llcccc}
\toprule
Pipeline & Metric & Uniform & Independent & Proportional & \SketchBudget{} \\
\midrule
P1: BF$\to$CMS & MAE $\downarrow$ & $1.48 \pm 0.01$ & $1.31 \pm 0.03$ & $0.71 \pm 0.01$ & $\mathbf{0.39 \pm 0.02}$ \\
P2: CMS$\to$T$\to$HLL & Card.\ err $\downarrow$ & $9.72 \pm 2.56$ & $3.33 \pm 1.35$ & $\mathbf{0.97 \pm 0.60}$ & $2.44 \pm 2.12$ \\
P3: BF$\to$CMS$\to$Sum & Abs.\ err $\downarrow$ & $1544 \pm 13$ & $1269 \pm 55$ & $678 \pm 32$ & $\mathbf{390 \pm 8}$ \\
\bottomrule
\end{tabular}
\end{table}

\SketchBudget{} achieves the lowest error on P1 and P3, reducing error by 45.2\% and 42.4\% respectively compared to the best baseline (Proportional). The improvement over Uniform is even larger: 73.6\% (P1) and 74.7\% (P3). These improvements are statistically significant ($p < 10^{-4}$, paired $t$-test).

On P2, however, the Proportional baseline achieves lower error ($0.97 \pm 0.60$) than \SketchBudget{} ($2.44 \pm 2.12$). We analyze this failure mode in Section~\ref{sec:p2_analysis}.

\paragraph{Network Trace Results.}
On synthetic network trace data (Zipfian $\alpha=1.1$), results are consistent: \SketchBudget{} reduces error by 51.3\% (P1) and 49.8\% (P3) over the best baseline, while achieving a modest 16.3\% improvement on P2 (where it does outperform Proportional on this distribution).

\subsection{NASA HTTP Access Log Evaluation}

To validate external applicability, we evaluate on real NASA HTTP server logs ($10^6$ requests from ${\sim}51$K distinct hosts, with 349 heavy hitters at threshold $T=196$). Table~\ref{tab:nasa} shows results at 500KB budget.

\begin{table}[t]
\centering
\caption{Results on NASA HTTP access logs at 500KB budget. Lower is better ($\downarrow$). Best in \textbf{bold}.}
\label{tab:nasa}
\begin{tabular}{llcccc}
\toprule
Pipeline & Metric & Uniform & Independent & Proportional & \SketchBudget{} \\
\midrule
P1 & MAE $\downarrow$ & $7.42$ & $8.08$ & $3.71$ & $\mathbf{2.10}$ \\
P2 & Card.\ err $\downarrow$ & $55.6$ & $\mathbf{5.94}$ & $12.0$ & $\mathbf{5.94}$ \\
P3 & Abs.\ err $\downarrow$ & $14279$ & $17165$ & $7858$ & $\mathbf{4118}$ \\
\bottomrule
\end{tabular}
\end{table}

\SketchBudget{} reduces error by 43.4\% (P1) and 47.6\% (P3) compared to the best baseline. On P2, \SketchBudget{} matches the Independent allocator and outperforms both Uniform and Proportional---a better result than on synthetic Zipfian data where Proportional dominated. The real web traffic has a steeper frequency distribution with well-separated heavy hitters (349 out of 51K hosts), which reduces near-threshold classification difficulty.

\textbf{Note on variance}: NASA results show near-zero standard deviation across seeds because all runs use the same fixed dataset. Random seeds only affect insertion order, and all three sketch types (Bloom filter, CMS, HLL) produce identical outputs regardless of insertion order. In contrast, synthetic experiments generate fresh random streams per seed, yielding meaningful variance. NASA results thus measure deterministic performance on a fixed workload rather than robustness to data variation.

\subsection{Bound Tightness Analysis}

We compare our composition bounds against naive bounds across pipeline types and Zipfian skew parameters. The tightness ratio (bound/observed error) measures how conservative the bound is---a ratio of 1.0 would be perfectly tight; ratios $\gg 1$ indicate conservative bounds.

\begin{table}[t]
\centering
\caption{Relative improvement of our bounds over naive bounds (percentage reduction in tightness ratio) across Zipfian skew parameters at 500KB budget. Note: both bounds remain conservative in absolute terms (see text); these percentages measure relative, not absolute, tightness.}
\label{tab:tightness}
\begin{tabular}{lccccc}
\toprule
 & $\alpha=0.5$ & $\alpha=0.8$ & $\alpha=1.0$ & $\alpha=1.2$ & $\alpha=1.5$ \\
\midrule
P1: BF$\to$CMS       & 4.5\% & 4.3\% & 3.7\% & 2.2\% & 0.6\% \\
P2: CMS$\to$T$\to$HLL & 99.5\% & 99.2\% & 99.4\% & 99.7\% & 99.8\% \\
P3: BF$\to$CMS$\to$Sum & 41.7\% & 40.9\% & 35.9\% & 19.9\% & 2.1\% \\
\bottomrule
\end{tabular}
\end{table}

The results in Table~\ref{tab:tightness} reveal distinct regimes:
\begin{itemize}
    \item \textbf{P1}: Tight bounds are modestly tighter (0.6--4.5\%), with the improvement decreasing as skew increases. At high skew, the Bloom filter's FPR becomes negligible because the CMS error dominates---the composition correction has less impact.
    \item \textbf{P2}: Tight bounds are orders of magnitude tighter ($>99\%$ improvement). The naive bound overestimates by treating all distinct items as potential false heavy hitters, while the distribution-aware bound correctly identifies that most items have frequencies far below the threshold.
    \item \textbf{P3}: Tight bounds show 2--42\% improvement, with larger gains at lower skew. When the distribution is more uniform (low $\alpha$), more items contribute to Bloom filter false positives, making the average-vs-worst-case CMS estimate distinction more impactful.
\end{itemize}

Importantly, both tight and naive bounds remain conservative (ratios $\gg 1$ in all cases), indicating room for further tightening, particularly through higher-order distribution-dependent terms.

\subsection{Memory Efficiency}

Figure~\ref{fig:error_vs_budget} shows error as a function of memory budget across all allocation strategies. \SketchBudget{} consistently achieves the lowest error across budget levels for P1 and P3, with the advantage most pronounced at medium budgets (50KB--500KB).

\begin{figure}[t]
\centering
\includegraphics[width=\linewidth]{figures/error_vs_budget.pdf}
\caption{End-to-end error vs.\ memory budget on Zipfian ($\alpha=1.0$) data. \SketchBudget{} (red) achieves lower error across budget levels for P1 and P3. For P2, Proportional and Independent allocators are competitive or better at several budget points. Note log-log scale.}
\label{fig:error_vs_budget}
\end{figure}

For quantitative memory savings, comparing the budget required to match Uniform allocation's error at 500KB: \SketchBudget{} requires 47\% less memory for P1 and 72\% less for P3. For P2, the computed memory savings of 47\% is misleading: it is measured at small budget levels where all methods produce very high absolute errors ($>$16K cardinality error against a true heavy hitter count of ${\sim}40$), making the ``savings'' practically irrelevant for real deployments. We report this transparently and note that memory savings for P2 are only meaningful at budgets $\geq$100KB where cardinality errors become small enough to be useful.

\subsection{P2 Analysis: Why Proportional Wins}
\label{sec:p2_analysis}

The Estimate-Threshold-Count pipeline (P2) presents a qualitatively different optimization landscape than P1 and P3. The key insight is that P2's error is dominated by the CMS's ability to correctly classify items near the threshold, which depends on having sufficient CMS width---not on the CMS/HLL memory split.

The Proportional allocator assigns 90\% to CMS and 10\% to HLL, which happens to be near-optimal because: (1)~HLL achieves good cardinality estimates even with minimal memory (512 bytes gives $<5\%$ relative error), so allocating more to HLL has diminishing returns; (2)~CMS accuracy near the threshold improves rapidly with width, and the Proportional split maximizes CMS allocation.

\SketchBudget{}'s optimizer sometimes over-allocates to HLL (e.g., 8KB HLL at 500KB budget) because the theoretical error model underestimates how quickly the false heavy hitter count drops with CMS width for concentrated frequency distributions. This is a known limitation of our Markov-bound-based false heavy hitter estimate, which does not capture the sharp threshold behavior of Zipfian distributions.

\subsection{Ablation: Memory Allocation Breakdown}

Figure~\ref{fig:allocation} shows how each strategy distributes memory.

\begin{figure}[t]
\centering
\includegraphics[width=0.8\linewidth]{figures/allocation_breakdown.pdf}
\caption{Memory allocation breakdown at 500KB for each pipeline and strategy. \SketchBudget{} allocates 94\% to CMS for P1 (vs.\ 50\% Uniform, 70\% Proportional), recognizing that CMS error dominates the composition.}
\label{fig:allocation}
\end{figure}

For P1, \SketchBudget{} allocates only 30KB (6\%) to the Bloom filter and 470KB (94\%) to the CMS, a stark contrast to Uniform (50:50) and Proportional (30:70). This reflects the composition-aware insight: the Bloom filter needs only enough memory to achieve a negligible FPR, after which all marginal memory is better spent reducing CMS error.

\subsection{Ablation: Distribution Sensitivity}

Figure~\ref{fig:distribution} shows how \SketchBudget{}'s advantage varies with Zipfian skew on P2.

\begin{figure}[t]
\centering
\includegraphics[width=0.7\linewidth]{figures/distribution_sensitivity.pdf}
\caption{P2 cardinality error vs.\ Zipfian skew parameter at 500KB budget. At high skew ($\alpha \geq 1.2$), all methods achieve near-zero error because heavy hitters are well-separated from the threshold. \SketchBudget{} achieves the largest relative advantage at $\alpha=1.2$, but Proportional is competitive or better at $\alpha \leq 1.0$.}
\label{fig:distribution}
\end{figure}

At low skew ($\alpha \leq 0.8$), many items cluster near the threshold, creating the hardest regime for heavy-hitter detection. At high skew ($\alpha \geq 1.5$), a few dominant items are well above the threshold, and all allocators achieve near-zero cardinality error. At $\alpha=0.8$, \SketchBudget{} achieves a 72\% reduction vs.\ Uniform (9.6 vs.\ 34.6), but the Proportional allocator is slightly better (8.4). \SketchBudget{} shows its largest advantage at $\alpha=1.2$ (92\% reduction vs.\ the best baseline), where heavy hitters are well-separated but not trivially so, and the composition-aware allocation correctly optimizes the CMS/HLL split.

\subsection{Ablation: Greedy vs.\ Exact Optimizer}

We compare the greedy marginal heuristic ($\Delta m = 100$ bytes) against the scipy optimizer.

\begin{table}[t]
\centering
\caption{Greedy vs.\ scipy optimizer agreement on error. Runtime shown for 500KB budget.}
\label{tab:greedy}
\begin{tabular}{lcccc}
\toprule
Pipeline & Budget & Greedy Error & Scipy Error & Agreement \\
\midrule
P1 & 10KB  & 303.2 & 300.5 & 99.1\% \\
P1 & 500KB & 0.387 & 0.358 & 92.0\% \\
P1 & 1MB   & 0.045 & 0.040 & 86.9\% \\
\midrule
P2 & 500KB & 3.31  & 4.83  & 68.5\% \\
P2 & 1MB   & 0.65  & 1.08  & 60.1\% \\
\midrule
P3 & 500KB & 416   & 404   & 97.0\% \\
P3 & 1MB   & 79    & 47    & 31.9\% \\
\bottomrule
\end{tabular}
\end{table}

The greedy heuristic performs well for P1 ($\geq 87\%$ agreement) but shows significant disagreements on P2 and P3, particularly at larger budgets. For P3, the greedy heuristic under-allocates to the Bloom filter because early iterations see small marginal returns from BF memory, while the BF's error reduction is non-linear and accelerates at higher memory levels. \textbf{At small budgets on P3, the greedy heuristic can be dramatically worse}: at 50KB, greedy produces error of 408,650 vs.\ the optimizer's 28,498---a $14\times$ discrepancy, because the greedy approach assigns almost all memory to CMS, missing the critical BF allocation that would eliminate most false positives.

This demonstrates that while greedy allocation is adequate for simple two-stage pipelines where the error landscape is relatively smooth, it can fail catastrophically on pipelines with non-convex error compositions. We recommend using the exact optimizer for production deployments.

\paragraph{Runtime.}
The scipy optimizer completes in $<7$ms for all configurations. The greedy heuristic is fast for P1 and P3 ($<1$ms), but takes 1.2--5.4 seconds for P2 at 500KB due to the distribution-dependent error evaluation at each marginal step. For runtime scaling, the scipy optimizer runs in $<40$ms for synthetic pipelines with 2--10 stages.

\subsection{Ablation: Pipeline Depth}

Our experiments test depth 2 (P1: BF$\to$CMS) and depth 3 (P2: CMS$\to$Threshold$\to$HLL). We observe that \SketchBudget{}'s advantage is substantial at both depths---45\% improvement for depth-2 P1 and 42\% for the two-stage P3---but we acknowledge that \textbf{we evaluated only 2--3 stage pipelines}. Deeper pipelines (4--5 stages) would require implementing additional pipeline types and conducting further validation. We conjecture that the advantage of composition-aware allocation grows with depth, as more stages create more opportunities for error propagation interactions, but leave systematic depth scaling analysis to future work.

% ============================================================
% 5. DISCUSSION
% ============================================================
\section{Discussion and Limitations}
\label{sec:discussion}

\paragraph{When does \SketchBudget{} help?}
Our results show a clear pattern: composition-aware allocation is most beneficial for pipelines where filtering errors (BF false positives) propagate multiplicatively to downstream estimation stages (P1, P3). In these cases, the optimizer recognizes that a small Bloom filter investment (6--10\% of budget) to eliminate false positives yields outsized reductions in CMS error. The benefit is less clear for threshold-based pipelines (P2), where error is dominated by the frequency distribution near the threshold rather than by inter-stage error propagation.

\paragraph{Limitations of the error model.}
Our composition bounds, while provably tighter than naive bounds, remain conservative in absolute terms. Tightness ratios range from 13--61$\times$ (P1) to 1537$\times$ (P2) at $\alpha=1.0$---the bounds dramatically overpredict the actual error. This conservatism arises from three sources: (1)~worst-case analysis over hash function families, which is inherent to the sketch error model; (2)~our linear approximation for CMS error probabilities, which does not capture the sharp exponential decay beyond the expected error; and (3)~for P2, the Markov-based false heavy hitter bound, which is very loose for concentrated frequency distributions. Practically useful error prediction would require distribution-dependent concentration inequalities (e.g., McDiarmid's inequality applied to the composed error) or empirical calibration. Nevertheless, the bounds' monotonicity properties suffice for the allocation optimizer to find good memory splits, which explains why \SketchBudget{} achieves strong empirical results despite the absolute looseness of the bounds.

\paragraph{P2 failure mode.}
\SketchBudget{} underperforms the Proportional baseline on P2 with Zipfian ($\alpha=1.0$) data. The root cause is that the Markov-style bound for false heavy hitters overestimates their count, causing the optimizer to over-invest in HLL precision at the expense of CMS width. A more refined error model for the near-threshold behavior of CMS could address this.

\paragraph{Greedy heuristic reliability.}
The greedy heuristic shows large discrepancies with the exact optimizer on P3 at small budgets (up to $14\times$ worse) and on P2 at medium budgets. This is because the greedy approach makes locally optimal but globally suboptimal decisions when the error function is non-convex in the allocation space. For production use, we recommend the exact optimizer, which runs in under 1ms.

\paragraph{Synthetic vs.\ real data.}
While we validate on real NASA HTTP logs and find consistent results for P1 and P3, broader evaluation on diverse real-world streams (network traffic, database workloads, IoT telemetry) would strengthen external validity.

\paragraph{Stream scale.}
Our experiments use $10^6$-item streams, whereas production systems process $10^9$+ items. While the theoretical analysis scales to arbitrary stream lengths, empirical validation at larger scales would increase confidence in the framework's practical applicability.

% ============================================================
% 6. CONCLUSION
% ============================================================
\section{Conclusion}
\label{sec:conclusion}

We presented \SketchBudget{}, a framework for composition-aware error analysis and memory allocation in heterogeneous sketch pipelines. Our composition bounds formalize how errors propagate through filter, threshold, and aggregation operations. The allocation optimizer, which solves the resulting constrained optimization in under 1ms, reduces end-to-end error by 42--51\% compared to the best baseline for filter-based pipelines (P1, P3) at 500KB budget, and achieves equivalent accuracy with 47--72\% less memory.

Our analysis also revealed important negative results: for threshold-based pipelines (P2), the Proportional baseline remains competitive because error is dominated by CMS accuracy near the decision boundary rather than by inter-stage error propagation. This suggests that composition-aware allocation is most valuable when upstream filtering errors have multiplicative effects on downstream stages.

Future work includes: (1)~substantially tighter bounds via distribution-dependent concentration inequalities (e.g., McDiarmid) and empirical calibration, addressing the 13--1537$\times$ overprediction gap; (2)~extending the framework to deeper pipelines (4+ stages) and more complex DAG topologies; (3)~online adaptive allocation that adjusts memory distribution as stream statistics evolve; and (4)~integration with learned sketch components~\citep{mitzenmacher2018model, dolera2023learning} whose error characteristics may differ from classical analysis.

% ============================================================
% REFERENCES
% ============================================================
\bibliographystyle{plainnat}

\begin{thebibliography}{12}

\bibitem[Bloom(1970)]{bloom1970space}
Burton~H. Bloom.
\newblock Space/time trade-offs in hash coding with allowable errors.
\newblock \emph{Communications of the ACM}, 13(7):422--426, 1970.

\bibitem[Cormode and Muthukrishnan(2005)]{Cormode2005CountMin}
Graham Cormode and S.~Muthukrishnan.
\newblock An improved data stream summary: The {Count-Min Sketch} and its applications.
\newblock \emph{Journal of Algorithms}, 55(1):58--75, 2005.

\bibitem[Dobra et~al.(2004)]{dobra2004sketch}
Alin Dobra, Minos Garofalakis, Johannes Gehrke, and Rajeev Rastogi.
\newblock Sketch-based multi-query processing over data streams.
\newblock In \emph{Advances in Database Technology---EDBT 2004}, pages 551--568. Springer, 2004.

\bibitem[Dobra et~al.(2009)]{dobra2009multi}
Alin Dobra, Minos Garofalakis, Johannes Gehrke, and Rajeev Rastogi.
\newblock Multi-query optimization for sketch-based estimation.
\newblock \emph{Information Systems}, 34(6):518--536, 2009.

\bibitem[Dolera et~al.(2023)]{dolera2023learning}
Emanuele Dolera, Stefano Favaro, and Stefano Peluchetti.
\newblock Learning-augmented {Count-Min Sketches} via {Bayesian} nonparametrics.
\newblock \emph{Journal of Machine Learning Research}, 24(12):1--60, 2023.

\bibitem[Flajolet et~al.(2007)]{flajolet2007hyperloglog}
Philippe Flajolet, \'{E}ric Fusy, Olivier Gandouet, and Fr\'{e}d\'{e}ric Meunier.
\newblock {HyperLogLog}: The analysis of a near-optimal cardinality estimation algorithm.
\newblock In \emph{2007 Conference on Analysis of Algorithms (AofA)}, pages 127--146. DMTCS, 2007.

\bibitem[Greenwald and Khanna(2001)]{greenwald2001space}
Michael Greenwald and Sanjeev Khanna.
\newblock Space-efficient online computation of quantile summaries.
\newblock In \emph{Proceedings of the 2001 ACM SIGMOD International Conference on Management of Data}, pages 58--66. ACM, 2001.

\bibitem[Kraska et~al.(2018)]{kraska2018case}
Tim Kraska, Alex Beutel, Ed~H. Chi, Jeffrey Dean, and Neoklis Polyzotis.
\newblock The case for learned index structures.
\newblock In \emph{Proceedings of the 2018 International Conference on Management of Data (SIGMOD)}, pages 489--504. ACM, 2018.

\bibitem[Mitzenmacher(2018)]{mitzenmacher2018model}
Michael Mitzenmacher.
\newblock A model for learned {Bloom} filters, and optimizing by sandwiching.
\newblock In \emph{Advances in Neural Information Processing Systems}, volume~31, pages 464--473. 2018.

\bibitem[{NASA}(1995)]{nasa_http}
{NASA Kennedy Space Center}.
\newblock {HTTP} access logs, {July} 1995.
\newblock Available at \url{https://ita.ee.lbl.gov/html/contrib/NASA-HTTP.html}, 1995.

\bibitem[Zhu et~al.(2025)]{zhu2025mnemosyne}
Zichen Zhu, Yanpeng Wei, Ju~Hyoung Mun, and Manos Athanassoulis.
\newblock Mnemosyne: Dynamic workload-aware {BF} tuning via accurate statistics in {LSM} trees.
\newblock \emph{Proceedings of the ACM on Management of Data (PACMMOD)}, 3(3), 2025.

\end{thebibliography}

% ============================================================
% APPENDIX
% ============================================================
\appendix
\section{Reproducibility}
\label{app:reproducibility}

\paragraph{Hardware.}
All experiments run on a Linux server with 2 CPU cores and 128GB RAM. No GPU is used. Total computation time for all experiments is approximately 4 hours.

\paragraph{Software.}
Python 3.12 with NumPy 1.26, SciPy 1.12, matplotlib 3.8, mmh3 4.0, xxhash 3.4.

\paragraph{Sketch Implementations.}
\begin{itemize}
    \item \textbf{Bloom filter}: Bit array with $k$ independent hash functions (mmh3 with different seeds). Memory = $\lceil m/8 \rceil$ bytes where $m$ is the number of bits. Optimal $k = (m/n) \ln 2$.
    \item \textbf{Count-Min Sketch}: $d \times w$ array of 4-byte unsigned integer counters. Default $d=5$ rows. Width $w = \lfloor m_\text{bytes} / (4 \times d) \rfloor$. Error guarantee: $\varepsilon = e/w$, $\delta = e^{-d}$.
    \item \textbf{HyperLogLog}: $2^p$ registers of 1 byte each. Precision $p$ chosen as the largest such that $2^p \leq m_\text{bytes}$. Standard error $\approx 1.04/\sqrt{2^p}$.
\end{itemize}

\paragraph{Hyperparameters.}
\begin{itemize}
    \item Stream length: $10^6$ items per experiment.
    \item Universe size: 100K (Zipfian, network trace), ${\sim}81$K (NASA).
    \item P1 query set: 5,000 positives + 5,000 negatives.
    \item P2 threshold: $10\times$ average frequency.
    \item P3 watchlist size: 1,000 items.
    \item Memory budgets: $\{10, 50, 100, 500, 1000\}$ KB.
    \item Random seeds: $\{42, 123, 456, 789, 1024\}$.
    \item Greedy heuristic step size: $\Delta m = 100$ bytes.
    \item Scipy optimizer: \texttt{minimize\_scalar} with \texttt{bounded} method for P1/P3; grid search over HLL precisions for P2.
\end{itemize}

\paragraph{Data Generation.}
Zipfian streams: items drawn i.i.d.\ from $P(X=i) \propto i^{-\alpha}$ for $i \in \{1, \ldots, 100000\}$. Network trace: Zipf($\alpha=1.1$) with $10^5$ ``IP addresses.'' NASA logs parsed with regex extracting the hostname field from each access log line.

\end{document}
